\chapter{Troubleshooting Layer 2}\label{ch:tshootl2}
\bp{2.4}

\begin{outcomes}
By the end of this chapter you will be able to:
\begin{tight}
  \item \textbf{Apply} a repeatable order of checks to a layer 2 fault instead of
        guessing.
  \item \textbf{Select} the right \cmd{show} command for a symptom, and read what
        it returns.
  \item \textbf{Use} \cmd{show logging}, ping, extended ping and capture output
        together to localise a fault.
  \item \textbf{Diagnose} the six faults that account for most switched-network
        problems.
\end{tight}
\end{outcomes}

\begin{prereq}
All of Part~II. This chapter assembles \chref{vlans} to \chref{discovery} into a
method.
\end{prereq}

\begin{keyterms}
symptom $\bullet$ hypothesis $\bullet$ scope $\bullet$ errdisable $\bullet$
show logging $\bullet$ severity $\bullet$ divide and conquer $\bullet$
bottom-up $\bullet$ baseline
\end{keyterms}

\begin{examnote}[This chapter is the exam's heaviest verb]
Topic \tp{2.4} --- ``troubleshoot basic Layer 2/Layer 3 connectivity and device
operations using show commands (including show logs), ping, extended ping, trace
route, and packet capture output'' --- is one of eight \emph{troubleshoot} topics.
Together with the eight \emph{configure} topics they are 55\% of the paper. The
method below is worth more marks than any single technology in this book.
\end{examnote}

\section{Scope first, then layers}

\begin{conceptbox}[Ask ``how many?'' before you ask ``why?'']
The single most useful question is how much is broken, because it tells you where
to look before you have run a command.

\begin{tight}
  \item \textbf{One device} $\rightarrow$ that port, that cable, that host's
        configuration.
  \item \textbf{One VLAN, everywhere} $\rightarrow$ the VLAN itself: does it
        exist, is it shut, is it allowed on the trunks, does its SVI have an
        address?
  \item \textbf{Everything on one switch} $\rightarrow$ that switch's uplink, or
        spanning tree.
  \item \textbf{Everything, everywhere} $\rightarrow$ a loop, the root bridge, or
        a routing or gateway failure.
  \item \textbf{Intermittent, under load} $\rightarrow$ physical layer, duplex, or
        a saturated link. Almost never a VLAN.
\end{tight}
\end{conceptbox}

\begin{conceptbox}[Then work bottom-up]
Once you know the scope, walk the layers from the bottom. Each step is one
command and rules out a whole class of cause.

\begin{tightnum}
  \item \textbf{Is the port up?} \cmd{show interfaces status}
  \item \textbf{Is it clean?} \cmd{show interfaces} --- errors, duplex, speed
  \item \textbf{Is it in the right VLAN?} \cmd{show interfaces switchport}
  \item \textbf{Does the switch see the device?} \cmd{show mac address-table}
  \item \textbf{Does the VLAN reach the next switch?}
        \cmd{show interfaces trunk}
  \item \textbf{Is spanning tree forwarding?} \cmd{show spanning-tree vlan}
  \item \textbf{Is there a gateway?} \cmd{show ip interface brief}
\end{tightnum}
Resist starting at step 7 because the user said ``I cannot reach the internet''.
\end{conceptbox}

\section{The first three commands}

\begin{verify}{SW1}{show interfaces status}
Port      Name               Status       Vlan       Duplex  Speed Type
Gi0/1     Uplink to SW2      connected    trunk      a-full a-1000 10/100/1000
Gi0/5     Desk 5             connected    10         a-full  a-100 10/100/1000
Gi0/6     Desk 6             notconnect   10           auto   auto 10/100/1000
Gi0/7     Desk 7             err-disabled 10           auto   auto 10/100/1000
Gi0/8     Desk 8             connected    15         a-full a-1000 10/100/1000
\end{verify}

One screen, four different faults visible. \cmd{Gi0/6} \cmd{notconnect} --- nothing
plugged in, or a dead cable. \cmd{Gi0/7} \cmd{err-disabled} --- a protection
feature fired; find out which. \cmd{Gi0/8} is in VLAN 15, which may or may not be
intended --- compare against the design.

\begin{verify}{SW1}{show logging | include LINK|LINEPROTO|ERR|SPANTREE}
%LINK-3-UPDOWN: Interface GigabitEthernet0/6, changed state to down
%SPANTREE-2-BLOCK_BPDUGUARD: Received BPDU on port Gi0/7 with BPDU Guard
   enabled. Disabling port.
%PM-4-ERR_DISABLE: bpduguard error detected on Gi0/7, putting Gi0/7 in
   err-disable state
%CDP-4-NATIVE_VLAN_MISMATCH: Native VLAN mismatch discovered on
   GigabitEthernet0/1 (999), with SW2 GigabitEthernet0/1 (1).
\end{verify}

\begin{alertbox}[Read the log before you read anything else]
The log frequently contains the answer in plain English, timestamped. Above, it
names the errdisable cause and a native VLAN mismatch nobody had reported --- a
fault that is silently bridging VLAN 999 and VLAN 1 between two switches.

Filter it. \cmd{show logging | include} with the keywords above turns hundreds of
lines into the handful that matter.
\end{alertbox}

\section{The six faults, and what each looks like}

\begin{longtable}{@{}L{3.6cm}L{5.0cm}L{6.4cm}@{}}
\toprule
\rowcolor{primary!15}
\textbf{Fault} & \textbf{Symptom} & \textbf{Command that proves it} \\
\midrule
\endhead
Port in the wrong VLAN & One user isolated; can reach some hosts, not others &
\cmd{show interfaces switchport} --- check Access Mode VLAN \\
\rowcolor{lightbg}
VLAN missing from a trunk & One VLAN fails between switches; others fine &
\cmd{show interfaces trunk} --- compare allowed lists at both ends \\
Native VLAN mismatch & Nothing obviously down; VLANs leak &
\cmd{show cdp neighbors detail}, or the log message above \\
\rowcolor{lightbg}
Errdisabled port & One port dead after somebody plugged something in &
\cmd{show interfaces status err-disabled} \\
Spanning tree blocking & A path that should work does not; redundant link idle &
\cmd{show spanning-tree vlan N} --- look for \cmd{Altn BLK} \\
\rowcolor{lightbg}
Duplex or speed mismatch & Slow, not down; errors climbing under load &
\cmd{show interfaces} --- collisions, late collisions, CRC \\
\bottomrule
\end{longtable}

\begin{worked}[Working a fault end to end]
\textbf{Symptom.} ``Nobody in accounts can reach the file server since this
morning. Sales are fine.''

\textbf{Scope.} One VLAN, multiple users, more than one switch. That rules out a
single cable and points at the VLAN or its path.

\begin{tightnum}
  \item \cmd{show interfaces status} on the access switch --- accounts ports are
        \cmd{connected} in VLAN 30. Not a port problem.
  \item \cmd{show mac address-table vlan 30} --- addresses are being learned
        locally. The devices are alive and the switch sees them.
  \item \cmd{show interfaces trunk} on the uplink --- VLAN 30 appears in
        ``allowed'' but \emph{not} in ``forwarding and not pruned''.
  \item \cmd{show spanning-tree vlan 30} --- the uplink is \cmd{Altn BLK}.
  \item \cmd{show logging | include SPANTREE} --- a topology change overnight; a
        second switch was added and won the root election for VLAN 30.
\end{tightnum}

\textbf{Cause.} Not a VLAN fault at all. The new switch became root for VLAN 30
only, which moved the tree and put the accounts uplink into blocking. Sales, in
VLAN 20, still had the old root and were unaffected --- which is exactly why only
one VLAN broke.

\textbf{Fix.} Set the intended root explicitly for all VLANs, add root guard on
the ports facing access switches, and confirm the uplink returns to
\cmd{Desg FWD}.
\end{worked}

\begin{pitfall}
\begin{tight}
  \item \textbf{Changing something before you have measured.} If you reconfigure
        before capturing the current state you can no longer tell whether you
        fixed it or moved it.
  \item \textbf{Changing more than one thing at once.} Then you do not know which
        change worked, and you cannot write it up.
  \item \textbf{Believing the reporter's diagnosis.} ``The switch is broken''
        usually means ``my thing does not work''. Take the symptom, discard the
        theory.
  \item \textbf{Ignoring what still works.} The working half of the network is
        evidence. Sales working and accounts failing localised the fault above
        faster than any command did.
  \item \textbf{Skipping the log.} It is timestamped and often explicit.
\end{tight}
\end{pitfall}

\begin{breakfix}[Two symptoms, one cause, and a misleading first look.]
Users report intermittent slowness. One user cannot connect at all. Both started
after a contractor visited the comms room.

\begin{verify}{SW1}{show interfaces status | exclude connected}
Port      Name               Status       Vlan       Duplex  Speed Type
Gi0/7     Desk 7             err-disabled 10           auto   auto 10/100/1000
\end{verify}

\begin{verify}{SW1}{show logging | include ERR\_DISABLE|MACFLAP}
%PM-4-ERR_DISABLE: bpduguard error detected on Gi0/7, putting Gi0/7 in
   err-disable state
%SW_MATM-4-MACFLAP_NOTIF: Host 0011.2233.4455 in vlan 10 is flapping
   between port Gi0/7 and port Gi0/12
\end{verify}

\textbf{Diagnosis.} Both symptoms are one event. The contractor plugged an
unmanaged switch into \cmd{Gi0/7} and, to ``make sure it worked'', also patched it
into \cmd{Gi0/12}. That created a loop: the MAC flap is the loop being seen from
two ports, and the intermittent slowness is the resulting flooding. BPDU guard
then errdisabled \cmd{Gi0/7}, which is the one port that is definitively dead ---
and which, by shutting, largely saved the network.

\textbf{Fix.} Remove the second patch lead. Confirm the flapping stops. Then
recover the port with \cmd{shutdown} / \cmd{no shutdown} --- \emph{after} the
cause is gone, not before, or it will simply errdisable again. Finally, note that
BPDU guard did its job: without it this would have been a full broadcast storm
rather than one dead port.
\end{breakfix}

\begin{lab}{Fix five faults against the clock}{Packet Tracer}
\textbf{Topology.} Three switches, one router, six PCs across three VLANs ---
built, working, and then broken by a partner.

\begin{tasks}
  \item Build the network and confirm it works fully. \textbf{Capture a baseline}:
        \cmd{show interfaces status}, \cmd{show interfaces trunk},
        \cmd{show spanning-tree vlan} and \cmd{show ip interface brief} from every
        device. You cannot troubleshoot without knowing what right looks like.
  \item Hand the network to a partner, who introduces \textbf{five} faults from
        the table above without telling you which.
  \item Find and fix them, writing down for each: the symptom, the scope, the
        command that proved it, and the fix. Give yourself twenty minutes.
  \item Swap roles.
  \item Compare notes: which fault took longest, and would starting with scope
        rather than commands have found it faster?
  \item \textbf{Extension.} Introduce a fault that produces \emph{no} log message
        --- a wrong access VLAN is a good one --- and discuss why the log is a
        first stop rather than a complete answer.
\end{tasks}
\end{lab}

\begin{checkpoint}
\begin{tightnum}
  \item Users on one VLAN across three switches lose connectivity; other VLANs are
        fine. What are the two most likely causes, and which command separates
        them?
  \item A port is errdisabled. You issue \cmd{no shutdown} and it errdisables
        again within seconds. What have you forgotten?
  \item Why does knowing what \emph{still works} narrow a fault faster than
        knowing what is broken?
\end{tightnum}
\end{checkpoint}

\begin{chaptersummary}
\begin{tight}
  \item Establish scope first --- one device, one VLAN, one switch, or everything.
        Scope points at the layer before you type anything.
  \item Then work bottom-up: port up, port clean, right VLAN, MAC learned, VLAN on
        the trunk, spanning tree forwarding, gateway present.
  \item Read \cmd{show logging} early. It is timestamped and often explicit.
  \item Six faults cover most cases: wrong VLAN, VLAN missing from a trunk, native
        VLAN mismatch, errdisable, spanning tree blocking, duplex mismatch.
  \item Capture a baseline while the network works. Change one thing at a time,
        and remove the cause before recovering an errdisabled port.
\end{tight}
\end{chaptersummary}

\begin{reviewq}
  \item \textbf{(MCQ)} A single VLAN fails between two switches while others work.
        Which command should you run first?
        (a)~\cmd{show ip route} (b)~\cmd{show interfaces trunk}
        (c)~\cmd{show version} (d)~\cmd{show mac address-table}
  \item \textbf{(MCQ)} Which status means a protection feature has shut a port?
        (a)~notconnect (b)~disabled (c)~err-disabled (d)~monitoring
  \item \textbf{(Short)} Explain why establishing scope before running commands
        makes troubleshooting faster, with an example.
  \item \textbf{(Short)} Give one fault from this chapter that produces a clear
        log message and one that produces none, and say how you would find the
        second.
  \item \textbf{(Applied)} Write a troubleshooting procedure of no more than eight
        steps that a first-line technician could follow for ``a user cannot reach
        the network'', naming the command at each step and the decision it
        resolves.
\end{reviewq}
